\subsection*{81. Use a Probationary Period to Improve Evidence, Not to Avoid Selection Discipline}

A lawful, clearly communicated probationary or initial engagement period can generate direct evidence of performance in the real environment. It should have realistic tasks, support, agreed standards, feedback, and defined decision points. It is not a license to avoid fair selection, pay below applicable requirements, or keep people perpetually insecure.

For a part-time reviewer, a paid and supervised 1st batch may provide evidence about turnaround, accuracy, learner-facing tone, and escalation. For an executive, a short provisional advisory project may demonstrate working relationship but cannot establish future leadership competence. High-stakes or irreversible decisions still require proper diligence before access to funds, private data, or public authority. Use the trial to learn rather than pretending it eliminates risk.

\subsection*{82. Reference Checks Should Target the Remaining Uncertainty}

A generic request for strengths and weaknesses often yields vague praise. After work samples and interviews, identify the precise claims worth corroborating:
\begin{itemize}[leftmargin=*]
\item \textbf{Instructor}: consistency of grading, treatment of struggling learners, reliability during live sessions.
\item \textbf{Senior engineer}: ownership of incidents, clarity of design decisions, influence without authority.
\item \textbf{Consultant}: scope discipline, truthfulness about results, contribution independent of the team.
\item \textbf{Sales leader}: quality of deals, forecast accuracy, respect for implementation limits.
\end{itemize}

Ask permission, describe the job, request observed incidents, and probe the reference's actual opportunity to observe the behavior. A thoughtful question is, ``What kind of assignment would you confidently trust this person to run without checking, and what would you still oversee?'' References are selected, partial, contextual data. Material negative information should be carefully examined, not automatically accepted or ignored.

\subsection*{83. Candidate Consent, Ownership, and Conflicts Are Genuine Hiring Variables}

A startup may be tempted to ask a candidate for source code from a previous employer, confidential client lists, unpublished research, training data, or private learner records to prove skill. Do not reward unethical disclosure. Provide neutral exercises; use open material or permissioned artifacts. In interviews, a candidate who refuses to share protected information may be demonstrating a value the firm should want.

Establish rights to work produced during any compensated selection project, keep evaluation materials limited to authorized purposes, and clarify ownership of future institutional intellectual property where lawful. For part-time and fractional roles, address availability, conflicts of interest, access permissions, confidentiality, and restrictions in other employment. The founder has similar obligations. No employment or consulting assessment should depend on misappropriated material.

\subsection*{84. AI-Assisted Assessments Need Explicit Tool Rules}

Candidates now routinely use AI to plan, write, code, summarize research, and prepare interviews. A rule of ``no AI'' may be appropriate for evaluating independent foundational competence in a specific task, but can be unrealistic for a job that actively requires AI-supported work. Conversely, permitting unrestricted tools may make it difficult to infer the candidate's reasoning.

Prefer a \textbf{2-mode protocol} when useful:
\begin{enumerate}[leftmargin=*]
\item \textbf{Mode 1}: observe unassisted reasoning on a small core problem.
\item \textbf{Mode 2}: allow tools as the job would, require disclosure of what was used, and test the candidate's ability to verify and critique the output.
\end{enumerate}
For a code task, ask why the generated patch is correct, what could break, and which tests catch it. For teaching, ask the candidate to diagnose why an AI-generated answer is pedagogically misleading. Avoid covert monitoring or speculative AI-detector scores as proof of dishonesty. State the rules before evaluation.

\subsection*{85. The Fairness Audit Must Follow the Whole Funnel}

Sourcing determines who learns that the opportunity exists; screening determines who is invited; assessment design determines whose abilities are visible; scheduling determines who can participate; interviews shape impressions; compensation and acceptance determine who can afford to join. A process can use structured questions and still be unfair if it systematically excludes qualified candidates through irrelevant requirements or inaccessible scheduling.

Audit the job relevance of degree requirements, spoken-language demands, equipment needs, geographic restrictions, unpaid tasks, and scoring criteria. Offer reasonable accommodations and protect privacy according to applicable law. When sample sizes permit, inspect pass rates and outcomes for potential disparities, using legal guidance appropriate to the jurisdiction. A small firm may lack data to make statistical claims, but it can still identify obvious design problems. Fairness is a constraint on \textbf{how evidence is obtained and used}, not an optional statement on the careers page.

\subsection*{86. Do Not Automate Judgments That Have Not Been Defined}

A vendor may offer automated résumé rankings, video-based personality scores, ``cultural similarity,'' or emotion inference. The founder should ask which job-related construct is measured, what validation evidence supports the specific use, how the tool performs across groups and conditions, what data is collected, how candidates can contest error, and whether an understandable alternative exists. A computer-generated score can be consistent and still predict the wrong thing.

For very small specialized hires, expensive automated screening is usually unnecessary. A structured evidence sheet and 2 calibrated evaluators may give more useful information. For larger operations, automation may help scheduling, administration, anonymization, and bounded checks, but \textbf{delegating ranking decisions} requires serious validity, fairness, and governance review. Tool sophistication should never substitute for judgment about what matters.

\subsection*{87. The 4 Hiring Outcomes, Including `Do Not Hire'}

A disciplined founder needs more than ``yes'' and ``no.'' A defensible decision may be:
\begin{enumerate}[leftmargin=*]
\item \textbf{Hire now}: when gates and economics are met.
\item \textbf{Contract or trial on bounded work}: when the need is real but long-term fit remains uncertain.
\item \textbf{Develop an internal person}: when potential and mentoring capacity are strong.
\item \textbf{Do not fill the role}: when no candidate meets essential requirements or the demand itself is unproven.
\end{enumerate}

The 4th decision is particularly difficult. Unfilled roles create inconvenience and delays, and founders dislike disappointing candidates. Yet a low-quality hire in a trust-intensive service can create months of repair. The willingness to remain understaffed temporarily (while simplifying the product or delaying demand) is sometimes the better business decision. This does not mean vacancies are free; their opportunity cost belongs in the decision model.

\subsection*{88. Make the Decision in This Order}

A practical final committee discussion should proceed through 9 disciplined steps:
\begin{enumerate}[leftmargin=*]
\item Verify genuine business need and affordable commitment.
\item Confirm noncompensatory gates.
\item Review direct job-work evidence.
\item Review structured behavioral and learning evidence.
\item Examine references, availability, and conflicts.
\item Read each evaluator's independent assessment.
\item Identify what remains uncertain and whether more evidence is worth obtaining.
\item Compare candidates and alternatives.
\item Write the decision rationale before making an offer.
\end{enumerate}

Do not start by asking everyone which candidate they ``liked best.'' That invites anchoring, status dominance, and retrospective justification. Begin with the rubric and job evidence, require dissenting assessors to explain the observations behind disagreement, and resist social pressure to reach consensus. A well-governed hiring decision can document uncertainty without pretending the evidence is complete.

\subsection*{89. Build a Post-Hire Feedback Loop Without Claiming Psychometric Proof}

After 30, 60, and 90 days, review the specific indicators predicted by the hiring process: review accuracy, task reliability, learner response, deployment incidents, sales qualification quality, operational defects, or manager team performance. Compare the hiring memo with observed outcomes. Record which signals were useful and which proved misleading, including the manager's contribution to success or failure.

With very few hires, there is insufficient data to estimate robust predictive validity coefficients or infer universal rules. Use outcomes for disciplined qualitative learning: were competencies misdefined, did the work change, was the assessment too narrow, did onboarding fail, or was the candidate inaccurately judged? Over years and larger samples, the organization may collaborate with qualified assessment professionals to evaluate selection processes more rigorously. The point is to make hiring \textbf{an improving institutional capability} rather than a series of isolated opinions.
